\documentclass[10pt, conference, a4paper, compsocconf]{IEEEtran}
\usepackage{amsmath,amsxtra,amssymb,amsthm,latexsym,amscd,amsfonts}
\usepackage[utf8]{vietnam}
\usepackage[top=2.7cm, bottom=5.4cm, left=2.72cm, right=3.0cm]{geometry}
\headsep=0.7cm
\headheight=1.0cm
\footskip=1.0cm
\voffset=-0.74cm

\usepackage{fancyhdr}
\pagestyle{fancy}
\renewcommand{\sectionmark}[1]{\markright{\MakeUppercase{#1}}{}}
\renewcommand{\headrulewidth}{0pt}

\ifCLASSINFOpdf
  \usepackage[pdftex]{graphicx}
  \usepackage{pgfplots}
  \pgfplotsset{compat=1.16}
  \DeclareGraphicsExtensions{.pdf,.jpeg,.png,.jpg}
\else
  \usepackage[dvips]{graphicx}
  \DeclareGraphicsExtensions{.eps}
\fi

\usepackage{array}
\usepackage{url}
\usepackage{float}
\usepackage[caption=false,font=footnotesize]{subfig}
\usepackage{tikz}
\usetikzlibrary{shapes.geometric,arrows.meta,positioning,fit,calc,backgrounds}
\usepackage{xcolor}
\usepackage{colortbl}

%-------- Mau sac --------
\definecolor{cBlue}{RGB}{173,216,230}
\definecolor{cOrange}{RGB}{255,200,100}
\definecolor{cGreen}{RGB}{144,238,144}
\definecolor{cPurple}{RGB}{221,160,221}
\definecolor{cYellow}{RGB}{255,239,150}
\definecolor{cCyan}{RGB}{175,238,238}
\definecolor{cSkyBlue}{RGB}{135,206,235}
\definecolor{cHdrBlue}{RGB}{70,130,180}
\definecolor{cHdrGreen}{RGB}{46,139,87}
\definecolor{cHdrOrange}{RGB}{210,105,30}
\definecolor{cRowA}{RGB}{220,235,250}
\definecolor{cRowB}{RGB}{235,255,235}

%-------- TikZ styles --------
\tikzset{
  nd/.style={
    rectangle, rounded corners=4pt,
    draw=gray!70, fill=#1,
    align=center, minimum height=0.58cm,
    font=\footnotesize
  },
  arr/.style={-Stealth, thick, draw=gray!60},
  arrt/.style={-Stealth, thin, draw=gray!50},
}

%-------- Header --------
\makeatletter
\def\ps@IEEEtitlepagestyle{%
  \def\@oddhead{\centering\mbox{\footnotesize\textit{%
    \;\;\;\;Bài tập lớn: Môn phát triển phần mềm mã nguồn mở
    -- Đại học Sài Gòn, 16/04/2026}}}%
  \def\@evenhead{\centering\mbox{\footnotesize\textit{%
    \;\;\;\;Bài tập lớn: Môn phát triển phần mềm mã nguồn mở
    -- Đại học Sài Gòn, 16/04/2026}}}%
  \def\@oddfoot{\scriptsize\thepage\hfil}%
  \def\@evenfoot{\scriptsize\hfil\thepage}%
}
\makeatother

%=================================================================
\begin{document}
\pagenumbering{arabic}
\fontsize{10}{12}\selectfont

\fancyhead[RE,LO]{\centering\footnotesize\textit{%
  Bài tập lớn: Môn Phát triển phần mềm mã nguồn mở
  -- Đại học Sài Gòn, 5-6/04/2026}}

%-------- Tieu de --------
\title{SmartDocAI: Xây Dựng Hệ Thống Hỏi Đáp Tài Liệu Thông Minh\\
  Tích Hợp RAG, Hybrid Search, GraphRAG và CoRAG với LLM Ollama}

\author{
  \IEEEauthorblockN{
    Nguyễn Gia Bảo\IEEEauthorrefmark{1},\;
    Nguyễn Hoàng Phúc\IEEEauthorrefmark{2},\;
    Trương Gia Hào\IEEEauthorrefmark{3},\;
    Phạm Tấn Khương\IEEEauthorrefmark{4},\;
    Đặng Huỳnh Minh Thái\IEEEauthorrefmark{5}
  }
  \IEEEauthorblockA{
    \IEEEauthorrefmark{1}3122560004\quad
    \IEEEauthorrefmark{2}3122560060\quad
    \IEEEauthorrefmark{3}3122410100\quad
    \IEEEauthorrefmark{4}3122410191\quad
    \IEEEauthorrefmark{5}3121410448\\
    \IEEEauthorblockA{GitHub: \IEEEauthorrefmark{1}https:/github.com/nguyengiabao91753, \IEEEauthorrefmark{2}https:/github.com/phucnguyen0704, \IEEEauthorrefmark{3}https:/github.com/svtankhuong}
      \IEEEauthorblockA{\IEEEauthorrefmark{4}https:/https://github.com/thaidang342003, \IEEEauthorrefmark{5}https://github.com/HaoWasabi}
    Khoa Công nghệ Thông tin -- Trường Đại học Sài Gòn
  }
}

\maketitle

%-------- Abstract --------
\begin{abstract}
Bài báo cáo trình bày toàn bộ quá trình nghiên cứu và phát triển hệ thống
\textbf{SmartDocAI} -- phần mềm hỏi đáp tài liệu thông minh được xây dựng
trong 10 tuần thuộc môn học Phát triển phần mềm mã nguồn mở. Hệ thống tích
hợp kiến trúc RAG (Retrieval-Augmented Generation) kết hợp mô hình ngôn ngữ
lớn Ollama chạy hoàn toàn cục bộ, cho phép người dùng đặt câu hỏi và nhận
câu trả lời chính xác dựa trên nội dung tài liệu PDF/DOCX được tải lên. Qua
ba giai đoạn phát triển liên tiếp, nhóm đã nghiên cứu và triển khai bốn
kiến trúc truy xuất: RAG cơ bản với FAISS và LangChain, Hybrid Search kết
hợp BM25, GraphRAG xây dựng đồ thị tri thức với thuật toán Leiden/Louvain
và BFS/DFS, và CoRAG với chiến lược truy xuất đa bước theo chuỗi suy luận.
Kết quả thực nghiệm cho thấy độ chính xác câu trả lời cải thiện từ 68\%
(RAG cơ bản) lên 91\% (CoRAG) trên tập câu hỏi kiểm tra.
\end{abstract}

\begin{IEEEkeywords}
RAG; LLM; Ollama; FAISS; BM25; LangChain; GraphRAG; CoRAG;
Hybrid Search; hỏi đáp tài liệu; mã nguồn mở
\end{IEEEkeywords}

\IEEEpeerreviewmaketitle

%=================================================================
\section{Giới thiệu và công trình liên quan}

Với sự bùng nổ của các mô hình ngôn ngữ lớn (LLM), nhu cầu xây dựng hệ
thống hỏi đáp thông minh trên tài liệu nội bộ ngày càng cấp thiết. Tuy
nhiên, LLM thuần túy gặp hai hạn chế lớn: kiến thức đóng băng tại thời
điểm huấn luyện và không có khả năng truy cập tài liệu đặc thù của tổ
chức. RAG (Retrieval-Augmented Generation)~\cite{lewis2020rag} được đề xuất
để giải quyết vấn đề này bằng cách kết hợp bước truy xuất tài liệu liên
quan trước khi sinh câu trả lời.

\textbf{SmartDocAI} hướng tới xây dựng hệ thống hỏi đáp tài liệu hoàn
chỉnh sử dụng toàn bộ công cụ mã nguồn mở: Ollama~\cite{ollama} chạy LLM
cục bộ không cần API trả phí, FAISS~\cite{johnson2019faiss} làm vector
database hiệu năng cao, và LangChain~\cite{langchain} để điều phối pipeline.

Các công trình liên quan nhóm tham khảo gồm: kiến trúc RAG chuẩn của
LangChain~\cite{langchain}; BM25~\cite{robertson2009bm25} -- thuật toán tìm
kiếm từ khóa thống kê; GraphRAG~\cite{edge2024graphrag} của Microsoft với
đồ thị tri thức; và CoRAG~\cite{wang2024corag} với chiến lược phân rã câu
hỏi và truy xuất nhiều vòng.

%=================================================================
\section{Thực nghiệm}

\subsection{Giai đoạn 1 (Tuần 1--6): Xây dựng nền tảng RAG}

Trong giai đoạn đầu, cả nhóm cùng nghiên cứu các thành phần cốt lõi và
xây dựng hệ thống RAG cơ bản.

\textbf{Xử lý tài liệu:} Hệ thống hỗ trợ hai định dạng phổ biến -- PDF
(dùng \texttt{PyMuPDF}) và DOCX (dùng \texttt{python-docx}). Văn bản được
chia thành các chunk với kích thước và overlap được điều chỉnh thực nghiệm.

\textbf{Embedding:} Các chunk được chuyển thành vector ngữ nghĩa bằng mô
hình \texttt{all-MiniLM-L6-v2} (384 chiều, mã nguồn mở). Độ tương đồng
giữa hai vector được đo bằng Cosine Similarity:
\begin{equation}
  \cos(\mathbf{u},\mathbf{v}) = \frac{\mathbf{u}\cdot\mathbf{v}}
    {\|\mathbf{u}\|\,\|\mathbf{v}\|}
\end{equation}

\textbf{FAISS -- nghiên cứu thuật toán:} Nhóm nghiên cứu ba loại chỉ mục
FAISS: \texttt{IndexFlatL2} (brute-force, chính xác 100\%),
\texttt{IndexIVFFlat} (phân cụm Voronoi để tăng tốc), và
\texttt{IndexHNSW} (đồ thị lớp phân cấp, hiệu năng cao nhất cho production).

\textbf{Ollama và Prompt Engineering:} Mô hình Llama-3 được triển khai cục
bộ qua Ollama. Nhóm tối ưu prompt theo template:
\textit{[Ngữ cảnh] + [Câu hỏi] + [Hướng dẫn trả lời]}.

\textbf{LangChain Pipeline:} Pipeline RAG chuẩn 5 bước:
Document Loading $\to$ Text Splitting $\to$ Embedding $\to$ Vector Storage
$\to$ Retrieval \& Generation (xem Hình~\ref{fig:rag-basic}).

\begin{figure}[H]
\centering
\begin{tikzpicture}[node distance=0.42cm]
  \node[nd=cBlue, text width=2.9cm]                (q)  {Câu hỏi người dùng};
  \node[nd=cOrange, text width=2.9cm, below=of q]  (em) {Embedding\\(MiniLM-L6-v2)};
  \node[nd=cCyan, text width=2.9cm, below=of em]   (re) {Truy xuất FAISS};
  \node[nd=cYellow, text width=2.9cm, below=of re] (tk) {Top-K Tài liệu};
  \node[nd=cPurple, text width=2.9cm, below=of tk] (ll) {LLM Ollama (Llama-3)};
  \node[nd=cGreen, text width=2.9cm, below=of ll]  (an) {Câu trả lời};

  \draw[arr] (q)--(em);
  \draw[arr] (em)--(re);
  \draw[arr] (re)--(tk);
  \draw[arr] (tk)--(ll);
  \draw[arr] (ll)--(an);

  \node[left=0.45cm of re, font=\tiny, text=gray!70, align=right]
    (doc) {Tài liệu\\PDF/DOCX};
  \draw[arrt] (doc)--(re);
\end{tikzpicture}
\caption{Pipeline RAG cơ bản với FAISS và Ollama -- Giai đoạn 1}
\label{fig:rag-basic}
\end{figure}

\subsection{Giai đoạn 2 (Tuần 6--8): Hybrid Search và kiểm soát Hallucination}

Vector search thuần túy gặp hạn chế với câu hỏi chứa tên chuyên ngành
hoặc mã số cụ thể. Giải pháp là tích hợp \textbf{BM25}, thuật toán tìm
kiếm từ khóa thống kê:
\begin{equation}
  \text{BM25}(D,Q) = \sum_{q_i \in Q} \text{IDF}(q_i)\cdot
  \frac{f(q_i,D)\,(k_1+1)}{f(q_i,D)+k_1\!\left(1-b+b\dfrac{|D|}{\overline{dl}}\right)}
\end{equation}

Kết quả từ FAISS và BM25 được hợp nhất qua
\textbf{Reciprocal Rank Fusion (RRF)}:
\begin{equation}
  \text{RRF}(d) = \sum_{r\in R}\frac{1}{k + r(d)}, \quad k=60
\end{equation}

Cơ chế kiểm soát hallucination theo kiến trúc \textbf{Self-RAG}: sau khi
LLM sinh câu trả lời, hệ thống tự kiểm tra độ trung thực với tài liệu
nguồn; nếu phát hiện nội dung không có căn cứ, hệ thống tự viết lại câu
hỏi và thực hiện lại bước truy xuất.

Bảng~\ref{tab:retrieval-compare} so sánh hiệu năng các phương pháp trên
tập 100 câu hỏi kiểm tra:

\begin{table}[H]
\renewcommand{\arraystretch}{1.3}
\caption{So sánh hiệu năng các phương pháp truy xuất}
\label{tab:retrieval-compare}
\centering
\small
\begin{tabular}{|>{\columncolor{cRowA}}l|c|c|c|}
\hline
\rowcolor{cHdrBlue!70}\color{white}\textbf{Phương pháp} &
  \color{white}\textbf{Recall@5} &
  \color{white}\textbf{MRR@10} &
  \color{white}\textbf{Tốc độ}\\
\hline
FAISS đơn thuần       & 0.71 & 0.65 & Nhanh \\
\hline
BM25 đơn thuần        & 0.68 & 0.62 & Nhanh \\
\hline
\rowcolor{cRowB}Hybrid (FAISS+BM25) & 0.84 & 0.78 & Trung bình \\
\hline
\rowcolor{cRowB}CoRAG & \textbf{0.91} & \textbf{0.86} & Chậm hơn \\
\hline
\end{tabular}
\end{table}

\subsection{Giai đoạn 3 (Tuần 8--10): GraphRAG và CoRAG}

\subsubsection{GraphRAG -- Đồ thị tri thức}

GraphRAG giữ nguyên FAISS+BM25 làm seed retrieval và bổ sung lớp
\textbf{Knowledge Graph} để mở rộng ngữ cảnh, phục vụ câu hỏi đòi hỏi
suy luận đa bước (multi-hop). Pipeline gồm 5 bước:

\begin{enumerate}
\item \textbf{Chunking + Embedding:} Tái sử dụng code từ giai đoạn 1.
\item \textbf{Entity \& Relation Extraction:} LLM/NER trích xuất thực thể
  và quan hệ từ mỗi chunk.
\item \textbf{Document Graph:} Node = chunk/entity; Edge = quan hệ tường
  minh hoặc độ tương đồng vượt ngưỡng.
\item \textbf{Community Detection:} Thuật toán Leiden/Louvain phát hiện
  cụm; LLM tóm tắt mỗi community.
\item \textbf{Graph Traversal:} BFS/DFS dọc đồ thị để thu thập ngữ cảnh
  đa bước, sau đó Rerank và đưa vào LLM.
\end{enumerate}

\begin{figure}[H]
\centering
\begin{tikzpicture}[node distance=0.42cm and 0.7cm]
  \node[nd=cBlue, text width=2.5cm]                 (q)  {Câu hỏi};
  \node[nd=cOrange, text width=2.5cm, below=of q]   (sd) {FAISS+BM25\\(Seed Retrieval)};
  \node[nd=cCyan, text width=2.5cm, below=of sd]    (ch) {Seed Chunks};
  \node[nd=cSkyBlue, text width=2.5cm, below=of ch] (rk) {Rerank};
  \node[nd=cPurple, text width=2.5cm, below=of rk]  (ll) {LLM (Ollama)};
  \node[nd=cGreen, text width=2.5cm, below=of ll]   (an) {Câu trả lời};

  \node[nd=cPurple!70, text width=2.5cm, right=1.0cm of ch]  (ge){Graph\\Expansion};
  \node[nd=cYellow, text width=2.5cm, below=0.42cm of ge]    (kg){Knowledge\\Graph};

  \draw[arr] (q)--(sd);
  \draw[arr] (sd)--(ch);
  \draw[arr] (ch)--(ge);
  \draw[arr] (ge)--(kg);
  \draw[arr] (kg.south west) to[bend right=18] (rk.north east);
  \draw[arr] (ch)--(rk);
  \draw[arr] (rk)--(ll);
  \draw[arr] (ll)--(an);

  
\end{tikzpicture}
\caption{Kiến trúc GraphRAG với lớp mở rộng Knowledge Graph}
\label{fig:graphrag}
\end{figure}

\subsubsection{CoRAG -- Truy xuất theo chuỗi suy luận}

Vấn đề cốt lõi của RAG truyền thống: nếu câu hỏi thiếu ý, bước retrieve
sẽ sai và toàn bộ pipeline thất bại. CoRAG giải quyết bằng cách phân rã
câu hỏi và thực hiện nhiều vòng retrieve độc lập:

\begin{itemize}
\item \textbf{Question Decomposition:} LLM tách câu hỏi phức tạp thành
  các sub-question, mỗi câu nhắm vào một khía cạnh cụ thể.
\item \textbf{Multi-round Retrieval:} Mỗi sub-question thực hiện retrieve
  riêng qua FAISS+BM25.
\item \textbf{Fusion \& Rerank:} Gộp tất cả tập tài liệu, loại trùng lặp
  và xếp hạng lại bằng RRF để chọn Top-N context tốt nhất.
\item \textbf{LLM Generation:} LLM đọc context đã fusion và trả lời toàn
  diện.
\end{itemize}

\begin{figure}[H]
\centering
\resizebox{\columnwidth}{!}{%
\begin{tikzpicture}[node distance=0.38cm and 0.6cm]
  \node[nd=cBlue, text width=3.0cm]              (q)  {Câu hỏi phức tạp};
  \node[nd=cOrange, text width=3.0cm, below=of q](dc) {Question Decomposition\\(LLM)};

  \node[nd=cYellow, text width=1.9cm,
    below left=0.5cm and 1.4cm of dc]  (s1){Sub-Q1};
  \node[nd=cYellow, text width=1.9cm,
    below=0.5cm of dc]                 (s2){Sub-Q2};
  \node[nd=cYellow, text width=1.9cm,
    below right=0.5cm and 1.4cm of dc] (s3){Sub-Q3};

  \node[nd=cCyan, text width=1.9cm, below=0.38cm of s1](r1){Retrieve\\FAISS+BM25};
  \node[nd=cCyan, text width=1.9cm, below=0.38cm of s2](r2){Retrieve\\FAISS+BM25};
  \node[nd=cCyan, text width=1.9cm, below=0.38cm of s3](r3){Retrieve\\FAISS+BM25};

  \node[nd=cSkyBlue, text width=1.9cm, below=0.38cm of r1](d1){Doc Set A};
  \node[nd=cSkyBlue, text width=1.9cm, below=0.38cm of r2](d2){Doc Set B};
  \node[nd=cSkyBlue, text width=1.9cm, below=0.38cm of r3](d3){Doc Set C};

  \node[nd=cPurple, text width=3.0cm, below=0.65cm of d2](fu){Fusion \& Rerank (RRF)};
  \node[nd=cPurple!55, text width=3.0cm, below=0.38cm of fu](ll){LLM (Ollama)};
  \node[nd=cGreen, text width=3.0cm, below=0.38cm of ll](an){Câu trả lời toàn diện};

  \draw[arr] (q)--(dc);
  \draw[arr] (dc)--(s1); \draw[arr] (dc)--(s2); \draw[arr] (dc)--(s3);
  \draw[arr] (s1)--(r1); \draw[arr] (s2)--(r2); \draw[arr] (s3)--(r3);
  \draw[arr] (r1)--(d1); \draw[arr] (r2)--(d2); \draw[arr] (r3)--(d3);
  \draw[arr] (d1)--(fu); \draw[arr] (d2)--(fu); \draw[arr] (d3)--(fu);
  \draw[arr] (fu)--(ll);
  \draw[arr] (ll)--(an);
\end{tikzpicture}
}
\caption{Kiến trúc CoRAG -- truy xuất đa bước theo chuỗi suy luận}
\label{fig:corag}
\end{figure}

\subsubsection{Tối ưu giao diện người dùng}

Nhóm bổ sung \textbf{Progress Bar} thời gian thực hiển thị từng bước xử
lý (Chunking $\to$ Embedding $\to$ Retrieving $\to$ Generating), giúp
người dùng theo dõi và hiểu rõ luồng hoạt động bên dưới hệ thống.

\subsubsection{So sánh chất lượng qua các giai đoạn}

Hình~\ref{fig:accuracy-chart} thể hiện độ chính xác câu trả lời đánh giá
thủ công trên 50 câu hỏi kiểm tra.

\begin{figure}[H]
\centering
\begin{tikzpicture}
\begin{axis}[
  width=\columnwidth, height=5.6cm,
  ybar, bar width=0.48cm,
  ylabel={Độ chính xác (\%)},
  ymin=55, ymax=100,
  symbolic x coords={RAG co ban, Hybrid Search, GraphRAG, CoRAG},
  xtick=data,
  xticklabels={RAG cơ bản, Hybrid Search, GraphRAG, CoRAG},
  x tick label style={font=\footnotesize, rotate=16, anchor=north east},
  nodes near coords,
  nodes near coords align={vertical},
  nodes near coords style={font=\scriptsize\bfseries},
  enlarge x limits=0.18,
  ymajorgrids=true, grid style={dashed, gray!40},
]
\addplot[fill=cBlue!85, draw=cHdrBlue] coordinates {
  (RAG co ban,    68)
  (Hybrid Search, 79)
  (GraphRAG,      85)
  (CoRAG,         91)
};
\end{axis}
\end{tikzpicture}
\caption{Độ chính xác câu trả lời theo từng phương pháp truy xuất}
\label{fig:accuracy-chart}
\end{figure}

Bảng~\ref{tab:full-compare} tổng hợp so sánh toàn diện bốn kiến trúc:

\begin{table}[H]
\renewcommand{\arraystretch}{1.3}
\caption{So sánh tổng quan bốn kiến trúc trong SmartDocAI}
\label{tab:full-compare}
\centering
\footnotesize
\setlength{\tabcolsep}{3.5pt}
\begin{tabular}{|>{\columncolor{cRowA}}p{1.7cm}|p{0.9cm}|p{1.0cm}|p{1.1cm}|p{1.0cm}|}
\hline
\rowcolor{cHdrBlue!70}
\color{white}\textbf{Tiêu chí} &
\color{white}\textbf{RAG} &
\color{white}\textbf{Hybrid} &
\color{white}\textbf{Graph-RAG} &
\color{white}\textbf{CoRAG}\\
\hline
Số lần retrieve     & 1 lần  & 1 lần  & Seed+Graph & N lần \\
\hline
Từ khóa chính xác  & Kém & \cellcolor{cRowB}Tốt & \cellcolor{cRowB}Tốt & \cellcolor{cRowB}Tốt \\
\hline
Suy luận đa bước   & Không & Không & \cellcolor{cRowB}Có & \cellcolor{cRowB}Có \\
\hline
Độ chính xác       & 68\% & 79\% & \cellcolor{cYellow!80}85\% & \cellcolor{cRowB}\textbf{91\%} \\
\hline
Chi phí runtime    & Thấp & Thấp & Cao & Trung bình \\
\hline
Giải thích được    & Thấp & Thấp & \cellcolor{cRowB}Cao & Trung bình \\
\hline
\end{tabular}
\end{table}

Hình~\ref{fig:time-chart} thể hiện phân bổ thời gian làm việc của nhóm:

\begin{figure}[H]
\centering
\begin{tikzpicture}
\begin{axis}[
  width=\columnwidth, height=5.0cm,
  ybar stacked, bar width=0.5cm,
  ylabel={Số giờ (ước tính)},
  symbolic x coords={GD1, GD2, GD3},
  xtick=data,
  xticklabels={Giai đoạn 1, Giai đoạn 2, Giai đoạn 3},
  x tick label style={font=\footnotesize},
  ymin=0, ymax=68,
  nodes near coords,
  nodes near coords style={font=\tiny},
  enlarge x limits=0.35,
  legend style={font=\footnotesize, at={(0.5,1.03)},
    anchor=south, legend columns=3},
  ymajorgrids=true, grid style={dashed, gray!40},
]
\addplot[fill=cBlue!85, draw=cHdrBlue]
  coordinates {(GD1,20)(GD2,12)(GD3,14)};
\addplot[fill=cOrange!80, draw=cHdrOrange]
  coordinates {(GD1,15)(GD2,10)(GD3,13)};
\addplot[fill=cGreen!70, draw=cHdrGreen]
  coordinates {(GD1,5)(GD2,6)(GD3,8)};
\addlegendentry{Nghiên cứu lý thuyết}
\addlegendentry{Lập trình}
\addlegendentry{Viết tài liệu}
\end{axis}
\end{tikzpicture}
\caption{Phân bổ thời gian làm việc nhóm theo từng giai đoạn (giờ)}
\label{fig:time-chart}
\end{figure}

%=================================================================
\section{Kết luận}

Qua 10 tuần phát triển, nhóm đã xây dựng thành công hệ thống SmartDocAI
với bốn kiến trúc truy xuất từ cơ bản đến nâng cao. Quá trình tiệm tiến:
RAG cơ bản giúp nhóm nắm vững nền tảng; Hybrid Search cải thiện độ bao
phủ từ khóa (68\% $\to$ 79\%); GraphRAG mở ra khả năng suy luận đa bước
qua đồ thị tri thức (85\%); CoRAG đạt độ chính xác cao nhất (91\%) bằng
chiến lược phân rã câu hỏi thông minh.

Điểm nổi bật là sử dụng hoàn toàn mã nguồn mở -- Ollama, FAISS, LangChain
-- cho phép triển khai không cần API trả phí, đảm bảo quyền riêng tư dữ
liệu và phù hợp với tổ chức vừa và nhỏ.

%=================================================================
\section{Hướng phát triển}

Nhóm dự kiến tiếp tục SmartDocAI theo các hướng sau:

\textbf{Kỹ thuật truy xuất:} Tích hợp Cross-encoder Reranker để cải thiện
xếp hạng tài liệu; bổ sung cơ chế khử trùng lặp chunk thông minh trong
bước Fusion; thử nghiệm mô hình embedding \texttt{vietnamese-sbert} (768
chiều, tối ưu cho tiếng Việt).

\textbf{Agentic RAG:} Nghiên cứu kiến trúc cho phép LLM tự quyết định
chiến lược retrieve phù hợp (RAG, CoRAG hay GraphRAG) tùy độ phức tạp của
câu hỏi đầu vào.

\textbf{Đánh giá tự động:} Tích hợp framework RAGAS để đo faithfulness,
answer relevancy và context precision tự động, thay thế đánh giá thủ công.

\textbf{Triển khai thực tế:} Đóng gói thành REST API để tích hợp vào ứng
dụng doanh nghiệp; xây dựng UI web hoàn chỉnh với lịch sử hội thoại và
quản lý tài liệu.

%=================================================================
\section{Tài liệu tham khảo}
\begin{thebibliography}{9}

\bibitem{lewis2020rag}
P.~Lewis et al., \emph{Retrieval-Augmented Generation for
Knowledge-Intensive NLP Tasks}, NeurIPS 2020.
\url{https://arxiv.org/abs/2005.11401}

\bibitem{johnson2019faiss}
J.~Johnson, M.~Douze, H.~Jégou, \emph{Billion-scale similarity search
with GPUs}, IEEE Transactions on Big Data, 2019.
\url{https://github.com/facebookresearch/faiss}

\bibitem{robertson2009bm25}
S.~Robertson và H.~Zaragoza, \emph{The Probabilistic Relevance Framework:
BM25 and Beyond}, Foundations and Trends in Information Retrieval, vol.~3,
2009.

\bibitem{langchain}
Harrison Chase et al., \emph{LangChain: Building applications with LLMs
through composability}, 2022.
\url{https://github.com/langchain-ai/langchain}

\bibitem{ollama}
Ollama Team, \emph{Run Large Language Models Locally}, 2023.
\url{https://ollama.ai}

\bibitem{edge2024graphrag}
D.~Edge et al., \emph{From Local to Global: A Graph RAG Approach to
Query-Focused Summarization}, Microsoft Research, 2024.
\url{https://arxiv.org/abs/2404.16130}

\bibitem{wang2024corag}
X.~Wang et al., \emph{Chain-of-Retrieval Augmented Generation}, 2024.
\url{https://arxiv.org/abs/2501.14342}

\end{thebibliography}

\end{document}